
\section{Text Elements}\fbeg{}\begin{frame}{\ft{Text Elements as Structured Objects}}


\begin{sftxt}[.93\textwidth]

I use the term \q{Sentence Object Model} to describe text encoding built around 
sentence objects, but where objects can encompass a spectrum of distinct types 
both at higher and lower scales.  Among sentence objects themselves, 
there is usually a straightforward structure --- where sentences 
follow one another in sequence, belonging to encompassing paragraphs --- 
but sometimes the structure is more complex.  For instance, 
footnotes may contain multiple sentences or even paragraphs, 
but are in some sense children of the sentences where their 
mark appears --- or, if it comes after a sentence, the 
sequence of sentences effectively diverges into main text 
and footnote content.  Similarly, a block quote or 
enumerated list often functions like a child of the sentence 
that introduces it, and sometimes the preceding 
sentence is grammatically incomplete and has the effect 
of merging with the content that follows it.  Also, 
special content such as quotations, citations, bibliographic entries, 
and entries/locators for print indexes have their 
own structural requirements.  As these examples 
suggest, modeling the full organization of textual content 
requires an ability to represent sometimes complex 
objects and their interrelationships.

\end{ftxt}


\end{frame}

